\documentclass[11pt,a4paper]{article}

% Packages
\usepackage[utf8]{inputenc}
\usepackage[T1]{fontenc}
\usepackage{amsmath}
\usepackage{graphicx}
\usepackage{hyperref}
\usepackage{geometry}
\usepackage[round]{natbib}
\usepackage{booktabs}
\usepackage{longtable}
\usepackage{array}
\usepackage{ragged2e}
\usepackage{caption}

\geometry{margin=1in}

% Title information
\title{Your Article Title}
\author{Your Name}
\date{\today}

\begin{document}

\maketitle

\begin{abstract}
Your abstract goes here. This is a brief summary of your article.
\end{abstract}

\section{Introduction}

This paper presents the methodology to classify neighbourhoods based on the observed mobility behaviours of their residents. Similar to existing geodemographic classifications \citep{wyszomierski2024neighbourhood}.

\section{Literature Review}

\section{Methods}

% ============================================================
% Variable Reference Table
% NTS-based Travel Behaviour and Demographic Classification
% ============================================================
% Requires in preamble:
%   \usepackage{booktabs}
%   \usepackage{longtable}
%   \usepackage{array}
%   \usepackage{ragged2e}
%   \usepackage{xcolor}
%   \usepackage{caption}
% ============================================================

\begin{longtable}{>{\ttfamily\small}p{4.8cm} >{\small}p{2.2cm} >{\RaggedRight\small}p{6.8cm}}

\caption{Variables used in the travel behaviour and demographic clustering analysis.
         All variables are sourced from the National Travel Survey (NTS) 2024.
         Diary-derived variables are computed from the 7-day travel diary;
         self-reported variables are from the interview questionnaire.}
\label{tab:variables} \\

\toprule
\normalfont\textbf{Variable} & \textbf{Type} & \textbf{Description} \\
\midrule
\endfirsthead

\multicolumn{3}{l}{\small\textit{Table \ref{tab:variables} continued from previous page}} \\[4pt]
\toprule
\normalfont\textbf{Variable} & \textbf{Type} & \textbf{Description} \\
\midrule
\endhead

\midrule
\multicolumn{3}{r}{\small\textit{Continued on next page}} \\
\endfoot

\bottomrule
\endlastfoot

% ------------------------------------------------------------------
\multicolumn{3}{l}{\textit{\textbf{A. Individual demographic variables}}} \\[3pt]
% ------------------------------------------------------------------

Sex\_B01ID
  & Binary
  & Sex of respondent (male / female). \\[4pt]

Age\_B01ID
  & Ordinal
  & Age of respondent in 21 banded categories (under 1 to 80+). \\[4pt]

MaritalStat\_B01ID
  & Categorical
  & Legal marital status. \\[4pt]

EthGroupTS\_B02ID
  & Binary
  & Ethnic group. Two categories: White; Non-White. \\[4pt]

COB\_B01ID
  & Categorical
  & Country of birth in 7 categories. \\[4pt]

EcoStat\_B02ID
  & Categorical
  & Economic activity status in 6 categories. \\[4pt]

NSSec\_B03ID
  & Ordinal
  & National Statistics Socio-Economic Classification (NS-SEC)
    at the 5-category high-level breakdown. \\[4pt]

EdAttn3\_B01ID
  & Ordinal
  & Highest educational qualification attained. \\[4pt]

Carer\_B01ID
  & Binary
  & Whether the respondent provides unpaid care for a family member
    or friend. \\[4pt]

% ------------------------------------------------------------------
\multicolumn{3}{l}{\textit{\textbf{B. Household demographic variables}}} \\[3pt]
% ------------------------------------------------------------------

HHoldStruct\_B02ID
  & Categorical
  & Household structure in 6 categories. \\[4pt]

HHoldNumAdults
  & Count
  & Number of adults (aged 18+) in the household. \\[4pt]

HHoldNumChildren
  & Count
  & Number of children (aged under 18) in the household. \\[4pt]

Ten1\_B02ID
  & Categorical
  & Tenure type in 3 categories. \\[4pt]

AddressType\_B01ID
  & Categorical
  & Dwelling type. \\[4pt]

NumCarVan\_B02ID
  & Ordinal
  & Number of cars or light vans available to the household
    in 3 banded categories. \\[4pt]

HHIncQDS2008\_B01ID
  & Ordinal
  & Equivalised household income quintile derived from the
    diary sample, using 2008 bandings. \\[4pt]

HHoldGOR\_B02ID
  & Categorical
  & Government Office Region of the household's residence. \\[4pt]

Settlement2021EW\_B02ID
  & Categorical
  & Urban--rural classification of residence based on the
    2021 Census in 3 categories: urban, rural, and other. \\[4pt]

% ------------------------------------------------------------------
\multicolumn{3}{l}{\textit{\textbf{C. Modal profile (diary-derived)}}} \\[3pt]
% ------------------------------------------------------------------

walk\_share
  & Continuous
  & Proportion of diary trips made on foot, ranging from 0 to 1. \\[4pt]

cycle\_share
  & Continuous
  & Proportion of diary trips made by pedal cycle, ranging from 0 to 1. \\[4pt]

car\_driver\_share
  & Continuous
  & Proportion of diary trips as a car or van driver,
    ranging from 0 to 1. \\[4pt]

car\_pass\_share
  & Continuous
  & Proportion of diary trips as a car or van passenger,
    ranging from 0 to 1. \\[4pt]

pt\_share
  & Continuous
  & Proportion of diary trips by public transport (bus, rail,
    underground, and taxi), ranging from 0 to 1. \\[4pt]

active\_share
  & Continuous
  & Proportion of diary trips by active modes (walking and cycling
    combined), ranging from 0 to 1. \\[4pt]

main\_mode
  & Categorical
  & The single mode group accounting for the highest share of
    diary trips. One of: walk, cycle, car, public transport,
    or other. \\[4pt]

mode\_entropy
  & Continuous
  & Shannon entropy of the distribution of trips across the five
    mode groups ($H = -\sum p_i \ln p_i$). Higher values indicate
    greater modal diversity. \\[4pt]

% ------------------------------------------------------------------
\multicolumn{3}{l}{\textit{\textbf{D. Activity profile (diary-derived)}}} \\[3pt]
% ------------------------------------------------------------------

commute\_share
  & Continuous
  & Proportion of diary trips made for commuting purposes. \\[4pt]

business\_share
  & Continuous
  & Proportion of diary trips made for employer's business. \\[4pt]

education\_share
  & Continuous
  & Proportion of diary trips for education or escorting
    to education. \\[4pt]

shopping\_share
  & Continuous
  & Proportion of diary trips for shopping. \\[4pt]

escort\_share
  & Continuous
  & Proportion of diary trips for escorting purposes
    other than education. \\[4pt]

personal\_share
  & Continuous
  & Proportion of diary trips for personal business
    (e.g.\ medical appointments, banking). \\[4pt]

leisure\_share
  & Continuous
  & Proportion of diary trips for leisure and recreational
    activities. \\[4pt]

dominant\_activity
  & Categorical
  & The single trip purpose accounting for the highest share
    of diary trips. \\[4pt]

activity\_entropy
  & Continuous
  & Shannon entropy of the distribution of trips across the eight
    purpose categories. Higher values indicate a more varied
    daily activity programme. \\[4pt]

% ------------------------------------------------------------------
\multicolumn{3}{l}{\textit{\textbf{E. Trip characteristics (diary-derived)}}} \\[3pt]
% ------------------------------------------------------------------

mean\_trips\_per\_day
  & Continuous
  & Average number of trips per day, computed as total diary trips
    divided by 7. \\[4pt]

mean\_stages\_per\_trip
  & Continuous
  & Average number of stages per trip across the diary week.
    A proxy for journey complexity and multimodality. \\[4pt]

mean\_trip\_dist\_miles
  & Continuous
  & Mean trip distance in miles, excluding short walk legs
    (\texttt{TripDisExSW}). \\[4pt]

total\_dist\_week\_miles
  & Continuous
  & Total distance travelled across the diary week in miles,
    including all trip legs. \\[4pt]

local\_trip\_share
  & Continuous
  & Proportion of diary trips under 2 miles. Captures the degree
    of hyperlocal travel behaviour. \\[4pt]

long\_trip\_share
  & Continuous
  & Proportion of diary trips over 25 miles. Captures the frequency
    of longer-distance travel. \\[4pt]

peak\_share
  & Continuous
  & Proportion of trips departing during AM peak (07:00--08:59)
    or PM peak (16:00--17:59). \\[4pt]

weekend\_share
  & Continuous
  & Proportion of diary trips made on a Saturday or Sunday. \\[4pt]

mean\_start\_time\_mins
  & Continuous
  & Mean trip departure time expressed in minutes past midnight.
    Captures habitual early or late travel patterns. \\[4pt]

mean\_trip\_time\_mins
  & Continuous
  & Mean total trip duration in minutes, including waiting
    and walking time at each end. \\[4pt]

mean\_daily\_travel\_time\_mins
  & Continuous
  & Mean time spent travelling per day in minutes, computed
    as total weekly travel time divided by 7. \\[4pt]

mean\_party\_size
  & Continuous
  & Mean number of people travelling together per trip,
    derived from stage-level party size. \\[4pt]

solo\_travel\_share
  & Continuous
  & Proportion of trips made alone (party size of one). \\[4pt]

% ------------------------------------------------------------------
\multicolumn{3}{l}{\textit{\textbf{F. Self-reported mode use frequencies (interview)}}} \\[3pt]
% ------------------------------------------------------------------

freq\_car
  & Ordinal
  & Self-reported frequency of private car use, recoded to a
    0--6 numeric scale (0 = never; 6 = three or more days
    per week). \\[4pt]

freq\_bus
  & Ordinal
  & Self-reported frequency of local bus use (0--6 scale). \\[4pt]

freq\_train
  & Ordinal
  & Self-reported frequency of surface rail use (0--6 scale). \\[4pt]

freq\_walk
  & Ordinal
  & Self-reported frequency of walking for 20 minutes or more
    without stopping (0--6 scale). \\[4pt]

freq\_cycle
  & Ordinal
  & Self-reported frequency of bicycle use (0--6 scale). \\[4pt]

freq\_taxi
  & Ordinal
  & Self-reported frequency of taxi or minicab use
    (0--6 scale). \\[4pt]

freq\_phv
  & Ordinal
  & Self-reported frequency of app-based private hire vehicle
    use (e.g.\ Uber), recoded to the same 0--6 scale. \\[4pt]

% ------------------------------------------------------------------
\multicolumn{3}{l}{\textit{\textbf{G. Work from home and digital substitution (interview)}}} \\[3pt]
% ------------------------------------------------------------------

wfh\_any
  & Binary
  & Whether the respondent worked from home on at least one day
    during the diary week (1 = yes; 0 = no). \\[4pt]

wfh\_frequency
  & Ordinal
  & Usual frequency of working from home, recoded to a 0--6
    numeric scale (0 = never; 6 = three or more times
    per week). \\[4pt]

\end{longtable}

\subsection{}

\section{Results}
Present your results.

\subsection{A Subsection}
You can add subsections as needed.

\section{Discussion}
Discuss your findings.

\section{Conclusion}
Your conclusions.

\bibliographystyle{apalike}
\bibliography{references.bib}

\end{document}